%% kapa: exam=1 round=26 year=2015 period=2 form=A subject=appraisal_law subject_ko=부동산관계법규 seq=12 number=12
%% source: Q-Net 감정평가사 기출문제 (cst00309) artlSeq=5138001 file=extracted/1008/5138001/제26회_감정평가사 1차_기출/제26회_감정평가사 1차_기출/2교시_감평_A형.pdf
%% conversion: kordoc 4.24.51 -> tools/gen1.py (auto); review=proofread-sonnet; flags=none
\begin{question}{12}
국토의 계획 및 이용에 관한 법령상 용도지역에 관한 설명으로 옳은 것은?
\begin{choices}
  \item 도시지역, 관리지역, 농림지역 또는 자연환경보전지역으로 용도가 지정되지 아니한 지역에 대해서는 용도지역의 건폐율과 용적률 규정을 적용할 때에 자연환경보전지역에 관한 규정을 적용한다.
  \item 도시지역은 주거지역, 상업지역, 공업지역, 자연환경보전지역으로 구분하여 지정한다.
  \item 자연환경보전지역은 녹지지역, 보전관리지역으로 구분하여 지정한다.
  \item 「항만법」에 따른 항만구역으로서 도시지역에 연접한 공유수면은 관리지역으로 결정ㆍ고시된 것으로 본다.
  \item 도시지역에 대해서는「도로법」제40조에 따른 접 도구역 규정이 적용된다.
\end{choices}
\end{question}
\answer{1}
